% SPDX-FileCopyrightText: Copyright (c) 2026 NVIDIA CORPORATION & AFFILIATES. All rights reserved.
% SPDX-License-Identifier: Apache-2.0
%
% ROLE: how coefficients are found: objective, search spaces and their conditioning, CMA-ES
%   settings, how noise is handled, restarts and initial points, selection, the equal-budget
%   rule, the drift check that gates the ladder, and what makes runs bit-exact across hosts.
%   A reader should be able to rerun lr-train from this section plus the appendix.
% CLAIMS: settled method descriptions only. Training curves, selected configurations and the
%   drift matrix are results (sec:results, sec:context).
% EVIDENCE: WT/benchmarks/learned_routing/learned_routing/{train.py,space.py,train_cli.py};
%   WT/benchmarks/learned_routing/remote/{p2orch.py,node_train.sh}; CR/runs/phase2/spaces/*.yaml;
%   CR/runs/phase2/MANIFEST.json; CR/runs/phase2/scripts/{flipcal.py,bc_fit.py};
%   CR/facts/gate.json full_plan; CR/facts/pilot.json; CR/facts/phase2_prep.json;
%   CR/facts/remote.json; CR/CONTRACT.md A1, A10, A11, A12, A15, A16, A17;
%   cma 4.5.0 in WT/.venv (cma/options_parameters.py).
% STATUS: final. Results of the runs are in sec:results; this section is method only.
% LABEL NOTE: train-side pilot numbers carry "preliminary (pilot, train subset)" in words; the
%   \prelim macro says "validation split", which would be false for them.

\section{Training method}
\label{sec:training}

% src: WT/benchmarks/learned_routing/learned_routing/train.py (module docstring); CR/facts/gate.json
%      full_plan.budget, objective_cells_repeats; CR/facts/pilot.json decisions
Training searches a policy's parameters directly against the replayed objective. No part of the
procedure scores a single routing decision: a decision changes the queues and caches that later
requests see, so its effect is visible only in aggregate over a replay. Each policy
configuration is therefore a black box that maps to one number, the objective below, and the
search is derivative-free. The spaces are small (7 free coefficients for M1 and 22 for M1-v2, two
to six parameters per baseline), and one tool, \code{lr-train}, tunes every policy, learned or
heuristic, with the same settings (\cref{alg:train}). This section describes phase~2; the pilot
used the same method on a 12-cell train subset, with 480 evaluations and two restarts per policy
(\cref{sec:pilot}, in the appendix).

\subsection{Objective}
\label{sec:train-objective}

% src: train.py (objective_value, pair_score, eval_set); CR/facts/gate.json
%      full_plan.objective_cells_repeats ("pre-registered, LR-01"); CR/facts/HEADLINE_TEST.json per_cell_statistic
A configuration $\pol$ is scored by replaying it on every train cell and averaging the clipped
log-ratio of windowed goodput (\cref{eq:clr}) against the default router on the same workload
replicates:
\begin{equation}
  \label{eq:objective}
  \Jobj(\pol) = \frac{1}{|\Cset{train}|}\sum_{\cell \in \Cset{train}}
    \frac{1}{\Krep}\sum_{\krep \in \Kset{\gen}} \clr_{\cell,\krep}(\pol,\piref),
  \qquad
  \Kset{\gen} = \{(\gen\Krep + i) \bmod \Kpool : 0 \le i < \Krep\},
\end{equation}
with $|\Cset{train}| = 34$, $\Krep = 2$ and $\Kpool = 8$. Each term is a paired comparison: the
candidate and $\piref$ replay the same cell on the same common-random-number (CRN) replicate with
the same policy seed $\krep+1$ (\cref{sec:aisim}). Before scoring, \code{lr-train} checks that
every returned record carries its task's cell, replicate and policy hash, and that candidate and
reference cover the same (cell, replicate) set; a mismatch stops the run rather than producing an
unpaired score. After its first replay, $\piref$'s record comes from the result cache.

The objective is episodic: one value per configuration and generation, aggregated over minutes of
simulated traffic per cell, with no per-decision credit assignment. Every cell weighs the same
whatever its absolute goodput. The clipped log-ratio was fixed before training and is the same
per-cell statistic the headline test uses (\cref{sec:headline}); the contract's original form, the
arithmetic mean of ratios, is reported but not optimized, because a few cells in which $\piref$
nearly collapses dominate it (\cref{sec:objective}).

% src: train.py Trainer.run ("worst = min(finite) - 1.0"), Trainer.score (reference error raises);
%      CR/facts/gate.json criteria.1_error_rate.lr_train_candidate_failures = 0 (4 runs x 480 evals)
A candidate whose replay fails on any cell receives the generation's lowest finite objective minus
one, so CMA-ES ranks it last; a failed reference replay stops the generation, which is retried on
the next call because errors are not cached. In the pilot no candidate failed in 4 runs of 480
evaluations (\emph{preliminary (pilot, train subset)}).
% src: 34 cells x 2 replicates x (16 candidates + reference) = 1,156 replays; re-evaluation
%      34 x 2 x (2 + 1) = 204 (arithmetic; matches CR/facts/phase2_prep.json early_progress
%      replays_fresh 1156 and tasks_requested 1360 at generation 1)
% src: CR/facts/pilot.json measurements.seconds_per_evaluation.full_train_local
%      replay_s_per_candidate_evaluation 700.895
One generation replays 16 candidates and $\piref$ on 34 cells $\times$ 2 replicates, \num{1156}
replays, plus 204 for the re-evaluation diagnostic below; a single candidate evaluation on the
full train split costs about 700.9 replay-seconds, as measured in the pilot.

\subsection{Search spaces}
\label{sec:space}

% src: WT/benchmarks/learned_routing/learned_routing/space.py (module docstring, Coord, cma_bounds);
%      every CR/runs/phase2/spaces/*.yaml param has bounds (checked 2026-10-03)
\paragraph{Internal coordinates.}
Each policy's \code{space.yaml} lists its free parameters with bounds and a linear or log scale,
its pinned parameters, and CMA-ES options. CMA-ES searches internal coordinates $\searchvec$: each
bounded parameter maps to $[0,1]$, linearly or in log space, and integer parameters are rounded
after decoding. pycma's boundary transform keeps every sample inside the box, so every candidate
decodes to a valid policy file. Every phase-2 parameter is bounded.

\paragraph{M1.}
% src: CR/runs/phase2/spaces/m1_learned_v1.yaml (header comment); CR/facts/pilot.json decisions (M1 space);
%      CR/literature/LESSONS.md LR-05 (Hansen tutorial pdf p.24: sigma grows under random selection)
At $\temp = 0$ the decision (\cref{eq:decision}) is unchanged when $\theta$ is multiplied by a
positive constant, a direction along which the objective is exactly flat and along which CMA-ES's
step size tends to grow \citep{hansen2016tutorial}. M1 therefore pins $\theta_0 = -1$, which fixes
the scale and the default router's sign and makes the start $\theta = -\basis{0}$ the default
router (\cref{sec:parity}). The temperature is fixed at $\temp = 0$ rather than searched: with the
scale pinned, $\temp > 0$ is a different, stochastic policy and is a separate variant
(\cref{sec:robustness}). The tie-break mode is \code{one\_draw}. Seven coefficients remain free.

% src: CR/facts/pilot.json feature_scaling.pooled_within_set_sd (geomean over the 12 pilot-subset
%      cells of each cell's mean within-set sd; CR/runs/pilot/screen/feature_scales.json "pooled");
%      ratios = sd(f0)/sd(f): 25.6, 17.7, 5.33, 80.8, 109.6, 21.2, 3.22 (also flip_rate.json s_j);
%      bounds from CR/runs/phase2/spaces/m1_learned_v1.yaml
Raw coefficients have no common scale, because the features' spreads within a candidate set
differ by two orders of magnitude (\cref{tab:m1-space}). We bound each free coefficient at about
four times its \emph{scale ratio} $\mathrm{sd}(x_{\cdot,0}) / \mathrm{sd}(x_{\cdot,f})$, the
coefficient at which feature $f$ varies across a candidate set as much as the anchor does; at its
bound, a feature's term spreads four times as widely as the anchor's. The standard deviations are
measured within candidate sets, as the geometric mean over the pilot's 12 train cells of each
cell's mean, on candidate tables reconstructed from $\piref$'s per-request rows, because replay
does not log candidate tables. The reconstruction is approximate (router load is rebuilt from
arrivals, first tokens and completions, and overlap is exact only for the chosen worker and for
earlier turns of the same session), so the ratios set scales, not policy inputs. The build
stage's placeholder bounds of $\pm 8$ could not have let \code{kv\_load\_frac} or
\code{active\_requests\_s} matter.

% src: space.py (bounded linear: u = (x - low)/(high - low)); sigma0 0.004 -> physical std
%      0.004 x 2 x bound = 0.032 x ratio (m1_learned_v1.yaml comment); LESSONS.md LR-05
%      (Hansen tutorial App. A, pdf p.29; ARS sec. 3.2, pdf p.7-8)
Because each bound is proportional to the feature's scale ratio, one internal unit moves every
coefficient by the same multiple of its ratio, so the initial search distribution is isotropic in
units of decision-relevant spread. This is the per-variable scaling \citet{hansen2016tutorial}
recommends when initial ranges differ by orders of magnitude, and the analogue of the state
normalization without which \citet{mania2018ars} could not train a linear policy on their
hardest task.

\begin{table}[tbp]
  \centering\small
  \caption{M1's search space and its three initial points. Within-set sd: geometric mean over the
    pilot's 12 train cells of the mean within-candidate-set standard deviation, on reconstructed
    tables. Ratio: $\mathrm{sd}(x_{\cdot,0})/\mathrm{sd}(x_{\cdot,f})$. The search box is
    $|\theta_f| \le$ bound. Start s1 is $\theta = -\basis{0}$ (all free coefficients 0); s2 and s3
    are described in \cref{sec:restarts}.}
  \label{tab:m1-space}
  % src: CR/facts/pilot.json feature_scaling.pooled_within_set_sd; CR/runs/pilot/screen/flip_rate.json s_j;
  %      CR/runs/phase2/spaces/m1_learned_v1.yaml bounds; CR/facts/phase2_prep.json m1_inits.s2.theta,
  %      m1_inits.s3.theta (s2 rounded to 3 significant digits)
  \begin{tabular}{@{}rlrrrrr@{}}
    \toprule
    $f$ & Feature & Within-set sd & Ratio & Bound & s2 & s3 \\
    \midrule
    0 & \code{default\_logit\_scaled}     & 6.607   & 1     & pinned & $-1$     & $-1$ \\
    1 & \code{overlap\_frac}              & 0.2585  & 25.6  & 100    & 63.9     & 76.8 \\
    2 & \code{new\_prefill\_tokens\_k}    & 0.3742  & 17.7  & 70     & 4.18     & 0 \\
    3 & \code{active\_prefill\_tokens\_k} & 1.239   & 5.33  & 21     & 1.28     & 0 \\
    4 & \code{kv\_load\_frac}             & 0.08173 & 80.8  & 320    & 12.5     & 0 \\
    5 & \code{active\_requests\_s}        & 0.06031 & 110   & 440    & $-138$   & 0 \\
    6 & \code{session\_affinity}          & 0.3111  & 21.2  & 85     & $-0.324$ & 42.4 \\
    7 & \code{isl\_x\_prefill\_load}      & 2.053   & 3.22  & 13     & $-0.196$ & 0 \\
    \bottomrule
  \end{tabular}
\end{table}

\paragraph{Baselines.}
% src: CR/facts/phase2_prep.json decisions[1..3], spaces.*.free/fixed/base_parameters
Each tuned baseline searches the parameters its catalog entry exposes (\cref{tab:spaces}),
starting from its shipped defaults. Three rules shape these spaces.
\begin{itemize}
  \item \textbf{Gauge pins.} A parameter that only rescales another is fixed. These are the
    affinity block size of \policy{ramjet} (512 tokens), the prefix weight of
    \policy{llm-d-precise-prefix} (2), and the peak prefill rate of
    \policy{llm-d-optimized-baseline} (\num{15928} tokens/s), whose product with
    \code{max\_ttft\_penalty\_ms} is the only quantity that policy identifies.
  \item \textbf{Structural parameters} keep their shipped values: \policy{lmetric}'s hot-spot
    detection on, \policy{ramjet}'s relative basis, \policy{sticky-session}'s map of
    \num{65536} sessions, \policy{llm-d-precise-prefix}'s empty class weights, and
    \policy{llm-d-optimized-baseline}'s throughput-based TTFT source.
  \item \textbf{Router knobs.} Of the baselines, only \policy{sticky-session} reads the default
    cost's router-level knobs (for first turns and fallbacks), so both sticky variants also search
    the five M0 knobs over M0's ranges. Sticky-session and M0 at the same knob values gave identical
    results on a session-free Mooncake cell, both at the start point and at the upper corner of the
    box. \code{router\_queue\_threshold} keeps its default for every policy (\cref{sec:ablations}).
\end{itemize}
% src: CR/runs/phase2/spaces/m0_default_tuned.yaml (comment, bounds, init); CR/facts/gate.json
%      full_plan.search_spaces.M0; CR/facts/pilot.json measurements.cma_progress_curves ("prefill_load_scale
%      moved to 17-25 of a 50 upper bound")
M0's space is the pilot's, with the prefill-load scale $\bpf$ widened to $[0.05, 500]$ on a log
scale because the pilot's tuned values (17 to 25) sat near the top of the old range $[0.05, 50]$.
The decode request weight $\breq$ starts at 0.5, the bottom of its log range $[0.5, 4096]$, which
stands in for the shipped value 0.

\paragraph{Initial step size.}
% src: CR/runs/phase2/scripts/flipcal.py (docstring: SIGMAS, SAMPLES 8, TARGET 0.175, TV distance,
%      stateful policies, clipping); CR/facts/phase2_prep.json flip_calibration (rule, cells = the 12
%      pilot-subset cells, target 0.175); CR/facts/DEVIATIONS.md 2026-10-03 phase-2 prep "sigma0 for tuned baselines"
Parameter units are not comparable across policies, so $\sigmazero$ is set in decision space.
For a step size $\sigma$, we draw internal points from $\mathcal{N}(\searchvec_0, \sigma^2 I)$
clipped to the box, decode them, and measure the share of routing decisions that differ from the
policy's own decisions at its start $\searchvec_0$, as the total-variation distance between the
two decision distributions. The decisions are one-step counterfactuals on the reconstructed
candidate tables of the pilot's 12 train cells, with 8 draws per $\sigma$ and cell; stateful policies run
their own state machine over each cell in arrival order. $\sigmazero$ is the $\sigma$,
interpolated on a grid from 0.03 to 0.35, at which this flip rate reaches 17.5\%, so every policy
starts with a comparable amount of exploration measured in decisions. The literature review's
suggested band of 10--30\% is itself a hypothesis (LR-05).
% src: CR/facts/pilot.json feature_scaling.m1_flip_frac_by_c (c = 0.03 -> 0.196, c = 0.05 -> 0.239;
%      sigma0 0.004 corresponds to c = 0.032); CR/runs/phase2/spaces/m1_learned_v1.yaml comment ("~20%")
For M1, $\sigmazero = 0.004$ internal units, a standard deviation of 0.032 scale ratios per
coefficient, flips about 20\% of decisions relative to $\theta = -\basis{0}$.
% src: CR/facts/phase2_prep.json spaces.<name>.flip_calibration.flip_at_sigma0_interpolated,
%      target_reached_by_0.3; flip_calibration.rule (key depth: "flips 0.69 / 0.18 already at sigma 0.03")
Five baselines cannot reach the target by $\sigma = 0.3$, the step-size cap below, and start at
0.3. Around their shipped defaults their decisions barely respond to their parameters:
\policy{lmetric}'s flip rate at $\sigma = 0.3$ is 0.06\%. For \policy{chwbl}'s and
\policy{dualmap}'s hashed prefix depths, any change re-draws the whole prefix-to-worker map (flip
rates of 0.69 and 0.18 already at $\sigma = 0.03$), so those two coordinates are held at their
defaults during the measurement and share their space's $\sigmazero$.

\begin{table}[tbp]
  \centering\footnotesize
  \caption{Search spaces of the tuned policies. Dim.: number of free coordinates.
    Flips: share of one-step decisions that change at $\sigmazero$ relative to the start, on
    reconstructed candidate tables; the calibration target is 0.175. $^{\dagger}$Hashed prefix
    depth held at its default during the measurement. Every space starts at the policy's shipped
    defaults, except M1, which has three starts (\cref{tab:m1-space}).}
  \label{tab:spaces}
  % src: CR/facts/phase2_prep.json spaces.<file>.{dimension,sigma0,free,fixed,base_parameters,
  %      flip_calibration.flip_at_sigma0_interpolated}; M1 flips from CR/facts/pilot.json
  %      feature_scaling.m1_flip_frac_by_c (measured against theta = -e0)
  \begin{tabularx}{\linewidth}{@{}lYrrr@{}}
    \toprule
    Policy & Free parameters (pinned) & Dim. & $\sigmazero$ & Flips \\
    \midrule
    M1 (\lc{}, \code{v1}) & $\theta_1,\dots,\theta_7$ ($\theta_0 = -1$, $\temp = 0$) & 7 & 0.004 & $\approx 0.20$ \\
    M0 (\defaultcf{}) & \code{overlap\_score\_credit}, \code{overlap\_score\_credit\_decay},
      \code{prefill\_load\_scale}, \code{decode\_active\_request\_weight}, \code{router\_temperature} & 5 & 0.11 & 0.169 \\
    \policy{sticky-session} (hard) & the five M0 knobs & 5 & 0.17 & 0.179 \\
    \policy{sticky-session} (bounded) & \code{load\_factor} and the five M0 knobs & 6 & 0.067 & 0.175 \\
    \policy{llm-d-precise-prefix} & \code{weights.queue}, \code{weights.kv\_cache\_utilization},
      \code{prefix\_match\_length\_weight}, \code{prefix\_match\_length\_scale\_tokens}
      (\code{weights.prefix} $= 2$) & 4 & 0.16 & 0.177 \\
    \policy{chwbl} & \code{prefix\_tokens}, \code{load\_factor} & 2 & 0.16 & 0.173$^{\dagger}$ \\
    \policy{ramjet} & \code{alpha}, \code{max\_affinity\_blocks}, \code{load\_unit\_tokens}
      (\code{affinity\_block\_tokens} $= 512$) & 3 & 0.30 & 0.133 \\
    \policy{dynamo-two-tier-cost-fn} & \code{cache\_threshold}, \code{balance\_abs\_threshold},
      \code{balance\_rel\_threshold} & 3 & 0.30 & 0.083 \\
    \policy{llm-d-optimized-baseline} & \code{affinity\_threshold}, \code{max\_ttft\_penalty\_ms},
      \code{queue\_threshold\_tokens} (peak prefill rate $= \num{15928}$) & 3 & 0.30 & 0.034 \\
    \policy{dualmap} & \code{hash\_prefix\_blocks}, \code{pending\_prefill\_token\_budget},
      \code{window\_requests} & 3 & 0.30 & 0.018$^{\dagger}$ \\
    \policy{lmetric} & \code{class\_prefix\_blocks}, \code{window\_requests} & 2 & 0.30 & 0.0006 \\
    \bottomrule
  \end{tabularx}
\end{table}

\paragraph{Step-size cap.}
% src: CR/runs/phase2/spaces/*.yaml (cma maxstd 0.3 in every baseline space; absent in the three M1
%      spaces); CR/facts/pilot.json measurements.cma_progress_curves ("seed 1's sigma then diverged
%      (0.18 -> 10.8 internal units ...): flat directions in decode_active_request_weight/temperature"),
%      recommendations_for_phase2[1]
% src: cma 4.5.0 cma/options_parameters.py:84-85 (maxstd None, maxstd_boundrange 1/3) and
%      cma/evolution_strategy.py:1179-1183, 1539-1564 (_stds_into_limits rescales sigma_vec, not sigma);
%      checked 2026-10-03: CMAEvolutionStrategy with bounds [0,1]^7 has opts['maxstd'] = 1/3 per coordinate
% check: whether the pilot's M0 seed-1 per-coordinate stds actually reached the 1/3 cap (history logs only sigma)
Every baseline space caps pycma's per-coordinate standard deviation (\code{maxstd}) at 0.3 internal
units. The pilot recommended the cap after M0's first run, whose global step size grew from 0.18 to
10.8 internal units after its validation score had plateaued, along nearly flat directions in the
decode request weight and the temperature (\emph{preliminary (pilot, train subset)}). M1's space,
whose step size stayed small in the pilot, sets no explicit cap. Two details qualify the measure.
pycma 4.5.0 already caps each bounded coordinate at one third of its range by default, so on these
$[0,1]$ coordinates the explicit cap tightens the default from $1/3$ to 0.3, and the default
applies to M1. The cap also limits the per-coordinate spread of the samples, not the global step
size $\sigma$, which can still drift.

\subsection{CMA-ES and noise handling}
\label{sec:noise}

% src: CR/facts/gate.json full_plan.budget; CR/runs/phase2/MANIFEST.json runs[*].args; train.py
%      _new_state (options seed, bounds, tolflatfitness 10, tolfun 0, tolfunhist 0, tolx 1e-6);
%      cma 4.5.0 defaults popsize 4 + 3 ln n, CMA_mu popsize // 2 (checked: mu = 8 at popsize 16;
%      4 + floor(3 ln n) = 6..9 for n = 2..7)
We use CMA-ES \citep{hansen2016tutorial} from the \code{cma} package, version 4.5.0, with
population size $\popsize = 16$, the package's default of 8 parents, and a budget of
$\budget = 400$ objective evaluations per run, which is 25 generations. For the 2 to 7 free
coordinates of these spaces the package's default population would be 6 to 9. Its function-value
stopping rules are switched off, so a run normally ends when the budget is spent; the equal-budget
check (\cref{sec:budget}) flags any finished run that did not spend exactly $\budget$.

The objective is noisy in a specific way. Replay is deterministic, so a configuration scored on a
given set of cells and replicates always gets the same value. What varies is which replicates a
generation draws (\cref{sec:aisim}) and, behind that, which workloads the train split happens to
contain. Four measures keep this noise from steering the search.
\begin{enumerate}
  \item \textbf{Common random numbers within a generation.} All 16 candidates of generation $\gen$
    and the reference $\piref$ replay the same cells and replicates $\Kset{\gen}$ with the same
    policy seeds, so CMA-ES ranks paired differences rather than independent draws. Scoring every
    candidate on the same fixed scenarios is the device of PEGASUS \citep{ng2000pegasus}. In a
    two-server load-balancing example, conditioning on the exogenous input sequence lowered
    policy-gradient variance 50-fold \citep{mao2019inputdriven}; for rank-based CMA-ES the analogue
    is paired scoring on shared replicates. The \code{one\_draw} tie-break keeps the random stream
    aligned across candidates whose decisions differ (\cref{sec:parity}).
    % src: CR/literature/LESSONS.md LR-03 (Mao ICLR'19 pdf p.3, p.5: 50x lower gradient variance;
    %      PEGASUS pdf p.3-7)
  \item \textbf{Rotating replicates.} $\Kset{\gen} = \{2\gen, 2\gen+1\} \bmod 8$: consecutive
    generations see different replicates and each pair recurs every fourth generation, so the
    CMA-ES distribution, which accumulates information over generations, is not fit to one draw.
  \item \textbf{Population size.} $\popsize = 16$ is about twice the default. In the analysis
    behind UH-CMA-ES, raising the population size together with the number of parents, as the
    package does, is preferable to re-evaluating each candidate several times when step sizes adapt
    properly \citep{hansen2009uh}.
    % src: CR/literature/LESSONS.md LR-03 (UH-CMA-ES sec. II, pdf p.4)
  \item \textbf{Rank-stability diagnostic.} After each update, the two best candidates are
    re-scored on the same cells with the next generation's replicates $\Kset{\gen+1}$, and the
    change in their order is logged. The re-scored share, $2/\popsize = 0.125$, equals
    UH-CMA-ES's default rate $\max(0.1, 2/\popsize)$ \citep{hansen2009uh}, but unlike UH-CMA-ES
    nothing reacts to the measurement: neither the step size nor the number of evaluations per
    candidate changes. It runs identically for every policy, does not count against $\budget$,
    and is skipped in a generation whose wall-clock chunk ends during it. In the pilot, the top two
    candidates swapped in 7 to 11 of 30 generations per run (\emph{preliminary (pilot, train
    subset)}).
    % src: train.py Trainer._reeval (count = max(2, ceil(reeval_frac x popsize)); fresh_ks =
    %      (g K + K + j) mod P; "diagnostic"); es.tell uses only the K = 2 training values;
    %      CR/runs/phase2/MANIFEST.json args --reeval-frac 0.125; p2orch.py budget_report comment;
    %      CR/literature/LESSONS.md LR-03 (UH-CMA-ES sec. IV, pdf p.8-9: r = max(0.1, 2/lambda));
    %      CR/facts/pilot.json runs.{m0-s1,m0-s2,m1-s1,m1-s2}.uh_cmaes_top2_swaps (8, 11, 7, 11 of 30)
\end{enumerate}

% src: train.py Trainer._validate (best_so_far + xfavorite on val cells, ks = range(3)), run (val_due
%      every --val-every 5 generations); CR/runs/phase2/MANIFEST.json args --val-every 5 --val-replicates 3
Every 5 generations, \code{lr-train} scores two candidates on all 14 validation cells with
replicates $\krep = 0,1,2$ ($\Kval = 3$): the best train sample so far and the mean of the CMA-ES
distribution. The mean is included because it is not itself chosen by a noisy train score. The
run keeps the higher-scoring validated candidate as its selection (\cref{sec:selection}).
\Cref{alg:train} summarizes one run.

\begin{algorithm}[t]
\caption{Tuning one policy (\code{lr-train}, one restart). Every policy, learned or baseline, runs
  this with the same $\budget$, $\popsize$, $\Krep$, cells, validation schedule and selection rule.}
\label{alg:train}
\begin{algorithmic}[1]
\Require space (bounds, pins, $\sigmazero$, cap), start $\searchvec_0$, seed, cells $\Cset{train}$
  and $\Cset{val}$
\State initialize CMA-ES at $\searchvec_0$ with step size $\sigmazero$ and population $\popsize$
\State $\textit{best} \gets \bot$; \quad $\textit{selected} \gets \bot$
\For{$\gen = 0, 1, \dots, \budget/\popsize - 1$}
  \State $\Kset{\gen} \gets \{(\gen\Krep + i) \bmod \Kpool : i < \Krep\}$
  \State sample candidates $\searchvec_1,\dots,\searchvec_{\popsize}$; checkpoint \Comment{\code{ask}}
  \State replay every candidate and $\piref$ on $\Cset{train} \times \Kset{\gen}$; check CRN coverage
  \State compute $\Jobj$ per candidate (\cref{eq:objective}); a failed one gets the worst finite $\Jobj$ $-1$
  \State update CMA-ES with fitness $-\Jobj$ \Comment{\code{tell}}
  \State $\textit{best} \gets$ the candidate with the highest $\Jobj$ seen so far
  \State re-score the top 2 on $\Cset{train} \times \Kset{\gen+1}$; log rank changes \Comment{diagnostic only}
  \If{$(\gen + 1) \bmod 5 = 0$}
    \State score \textit{best} and the CMA-ES mean on $\Cset{val} \times \{0,1,2\}$
    \State $\textit{selected} \gets$ the validated candidate with the highest score so far
  \EndIf
  \State checkpoint (CMA-ES state and random state)
\EndFor
\State \Return \textit{selected}, and \textit{best} for reference
\end{algorithmic}
\end{algorithm}

\subsection{Restarts and initial points}
\label{sec:restarts}

% src: CR/facts/phase2_prep.json decisions[0], m1_inits.restart_to_wave; CR/runs/phase2/MANIFEST.json
%      (seed = restart, wave = restart for tier A); CR/facts/gate.json full_plan.scheduling
Each tuned policy runs $\restarts = 3$ restarts with CMA-ES seeds 1, 2 and 3. Restart 1 of every
policy runs in the first wave, restart 2 in the second and restart 3 in the third, so a compute cut
at any point removes the same restart from every policy. Baselines start every restart at their shipped defaults; only the seed
differs. M1's three restarts start from three different points, all computed from train-split
data only (Amendment A11.1; \cref{tab:m1-space}).
\begin{description}
  \item[s1] $\theta = -\basis{0}$, the default router.
  \item[s2] A behavior clone of \policy{llm-d-precise-prefix}, the best of the untuned heuristics
    evaluated on validation in the pilot. It is the conditional-logit maximum-likelihood fit to that
    heuristic's choices: with $\Pr[i \mid \cand] \propto \exp(\theta^{\top}\fvec_i)$ over
    reconstructed candidate tables, it maximizes the mean over cells of the mean log-probability of
    the worker the heuristic chose, with $\theta_0$ pinned at $-1$ and $\theta_1,\dots,\theta_7$ held
    inside M1's box. The problem is convex and bounded and is solved by L-BFGS-B. The data are the
    heuristic's own replays at shipped defaults on all 34 train cells, replicate 0: \num{144048}
    decisions in a stride sample of each cell. The fit chooses the heuristic's worker on 0.875 of
    these decisions (0.866 and 0.883 in a two-fold split by cell), against 0.512 for the default
    rule $\theta = -\basis{0}$, and no coefficient sits at its bound.
    % src: CR/runs/phase2/scripts/bc_fit.py (docstring: model, pinning, bounds, L-BFGS-B, equal cell
    %      weights, stride sampling); CR/facts/phase2_prep.json m1_inits.s2_fit (decisions 144048, cells 34,
    %      agreement_in_sample 0.8748, agreement_cv_2fold_by_cell 0.8659/0.8827, agreement_theta0_default_rule
    %      0.5121, at_bound []); CR/facts/pilot.json mde.learned_vs_val_best_default_arm_heuristic.comparator
  \item[s3] A cache-heavy point: $\theta_1 = 3 \times 25.6 = 76.8$ and $\theta_6 = 2 \times 21.2 =
    42.4$, so that across a candidate set the overlap term spreads three times and the
    session-affinity term twice as widely as the anchor, with all other free coefficients 0, so
    load enters only through the anchor. It targets the agent workloads on which the pilot's M1
    trailed the cache-aware heuristic (by $-0.019$ in mean clipped log-ratio on AgentX, \prelim).
    It was preferred to a fit to the second-best heuristic family because \policy{lmetric}'s
    product of uncached tokens and load is not representable with \code{v1} features, and a fit to
    \policy{ramjet} would be another linear cache-minus-load score close to s2.
    % src: CR/facts/phase2_prep.json m1_inits.s3 (theta, why), decisions[7]; CR/facts/pilot.json
    %      validation_vs_best_heuristic.llm-d-precise-prefix.M1_tuned.by_family_mean_clipped_log_ratio.agentx
    %      (-0.01898); CR/literature/LESSONS.md LR-07 action 3
\end{description}

% src: CR/facts/phase2_prep.json m1_inits.{s2,s3}.one_step_flips_vs_theta0 (0.4645, 0.4198),
%      s2_vs_s3_one_step_disagreement (0.1096), s3.agreement_with_precise_prefix (0.8722),
%      s2_fit.caveat, s2_fit.sensitivity_hash_aware_overlap (agreement 0.7847, theta1 40.3),
%      s2_agreement_by_family (agentx 0.797, mooncake 0.914, sessions 0.865)
On one-step decisions, s2 and s3 differ from the default rule on 46\% and 42\% of requests and from
each other on 11\%. The clone's agreement should be read with the reconstruction's bias in mind:
overlap is exact only for the chosen worker, which favors overlap-heavy rules. With a hash-aware
reconstruction the fit agrees on 0.785 of decisions and its overlap coefficient is 40 instead of
64, and the hand-set point s3 agrees with the heuristic on 0.872, almost as often as the fit.
Agreement is lowest on AgentX (0.797, against 0.914 on Mooncake and 0.865 on sessions).

% src: CR/CONTRACT.md A11.1 ("Not extra budget"), A12; CR/facts/phase2_prep.json m1_inits.fairness,
%      s2_fit.data ("22 fresh replays + 12 cached, 0 errors"); CR/literature/LESSONS.md LR-07
%      (Park sec. 3.1, pdf p.4); LR-15 (Seshadri et al. sec. 4, pdf p.6)
The initial points are not objective evaluations. The clone uses one replicate of the heuristic's
own train replays (22 fresh and 12 cached), not evaluations of M1, so every M1 restart still spends
exactly $\budget$. Bootstrapping a policy search from existing policies is a known lever
\citep{mao2019park}, as is initializing a richer choice model at a simpler model's maximum-likelihood
fit \citep{seshadri2019cdm}. The starts still give M1 information that a baseline's restarts do
not have, namely another policy's decisions. The operator approved them with disclosure
(Amendment A12): the test also reports \emph{M1-default-init}, M1 selected over restart s1 alone
(\cref{sec:secondary}).

\paragraph{Other learned arms.}
% src: CR/runs/phase2/MANIFEST.json (m2r2-*, m1cont-* tier B: warm start (theta*_M1, p = 0), context
%      CMA_stds about 0.3x theta's; budget 400, 3 seeds); CR/CONTRACT.md A15 (M1-noaff: theta_6 pinned 0,
%      same multi-start inits, 3 restarts at B = 400, post-hoc theta_6 = 0 check on val k0-2), A16.2 (AIS
%      league: 3 restarts at B = 400, selected on val k0-2), A17.1 (M1-v2 inits)
\Cref{tab:ladder} lists the other learned arms. Each is trained like M1, with $\budget = 400$,
three restarts, the same cells and the same selection rule; they differ in their spaces and starting
points, and all of them ran to completion.
\begin{itemize}
  \item \textbf{M2 rank 2 and M1 continued.} M2 starts at the selected M1 coefficients with $p = 0$
    and context step sizes 0.3 times $\theta$'s; its two named sources, \code{isl\_k} and
    \code{mean\_active\_prefill\_tokens\_k}, were fixed before any M2 evaluation, and $\theta_7$ is
    pinned at 0 because the (\code{isl\_k}, feature 3) context entry duplicates it. It is compared
    with M1 continued from its selected point for the same added 400 evaluations (LR-15).
  \item \textbf{M1-noaff.} M1's three starts with $\theta_6$ pinned at 0; a further check, without
    retraining, zeroes $\theta_6$ in the selected M1 (Amendment A15).
  \item \textbf{M1-v2.} Coefficients 0 to 7 keep M1's bounds, so M1-v2 nests M1; coefficients 8 to
    22 are bounded at four times their scale ratio, measured on candidate tables reconstructed from
    the default router's train replays. $\sigmazero = 0.0028$ flips 19.6\% of decisions at the s1
    start, the rate M1's space has at its own $\sigmazero$. The starts are s1, $\theta = -\basis{0}$;
    s2, the behavior clone of \policy{llm-d-precise-prefix} refit on \code{v2} features (with
    \code{hash\_home} pinned at 0 in the fit, because its offline reconstruction is approximate);
    and s3, \policy{lmetric}'s representable point
    $\theta = -\basis{0} - 20.7\,(\basis{\text{\code{log\_ptok}}} + \basis{\text{\code{log\_bs}}})$,
    90\% of the bound on those coordinates, which agrees with the pure LMetric rule on 0.924 of
    reconstructed decisions (Amendment A17.1).
  \item \textbf{M1 with a queue threshold} (ablation a). M1 continued's space plus
    \code{router\_queue\_threshold} on a log scale from 0.25 to 256, starting at 32, a value that
    replays identically to the unset knob on every train cell; its baseline counterpart tunes the
    same knob on top of the best baseline's selected configuration.
  \item \textbf{\ais{} league.} M1-ais and M2-ais on feature set \code{v3}, M0 with the \ais{}
    prefill-load model, and \policy{llm-d-optimized-baseline} with modeled TTFT are trained the same
    way on their own build and compete only in a secondary leaderboard (Amendment A16,
    \cref{sec:res-ais}).
\end{itemize}
% src: CR/facts/tierBC.json decisions.m2_rank2 (sources, form, warm_start), m1_continue, m1_noaff,
%      m1_v2 (bounds, sigma0 0.0028 = flip rate 0.196, inits; s3 "-e0 - 20.7 (e_log_ptok + e_log_bs),
%      90% of bound_8 (agreement 0.924 with the pure LMetric rule)"), ablation_a (q log [0.25, 256],
%      start 32, identity check), ais_league; counts "done 30"

\subsection{Sign constraints on load coefficients}
\label{sec:sign-constraint}

% src: CR/audits/phase2-mid/gaming-concentration.md F1; CR/audits/phase2-mid/fix.md ("Constraint",
%      "How (a clamp, not a box bound)", "Inits", "Equal footing"); CR/facts/tierBC.json decisions
%      (m1_v2 sign_rule, m2_rank2 sign_rule, ais_league sign_rule); CR/CONTRACT.md A11.2
The midpoint audit found that the first tuned M1 gamed the objective by attracting long requests
to a busy worker (\cref{sec:gaming}), which triggered a remedy registered in advance (Amendment
A11.2): re-run the affected restarts with load coefficients constrained so that load cannot
attract, at the same budget. With $\theta_0$ pinned at $-1$, the minimal such set for \code{v1} is
$\theta_3, \theta_4, \theta_5, \theta_7 \le 0$ (active prefill, KV load, active requests and the
prompt-length-by-prefill interaction), so that no load signal adds utility for any prompt length.
\begin{itemize}
  \item \textbf{A clamp, not a box bound.} The restarts s1 and s3 start with those coefficients at
    0, which on a box bound of 0 would sit on the edge of pycma's boundary transform, where its
    slope is zero and the first generation's steps on those coordinates shrink by about two orders
    of magnitude. The space therefore keeps the original box and clamps the decoded value,
    $x = \min(x, 0)$. CMA-ES searches the identical internal space, start, step size and seed as the
    unconstrained run, and generation 0 asked the same internal points; only the decoded policies
    differ.
  \item \textbf{Starts.} s1 and s3 already satisfy the constraint. The behavior clone s2 did not
    (its $\theta_3$ and $\theta_4$ were positive, so decode load attracted above about 24K tokens),
    so it was projected onto the constraint; its agreement with the cloned heuristic stayed at 0.879.
  \item \textbf{Later arms.} Every later learned arm uses the same rule: M1-v2 clamps every
    load-ordered feature (3, 4, 5, 7, 8, 10 and 12 to 21), M2 also clamps the load columns of its
    context directions so that every effective load weight stays non-positive, and the \ais{} league
    clamps its load and backlog features.
  \item \textbf{Equal footing.} The re-runs used the same seeds, budget, cells, replicates,
    validation schedule and objective; the unconstrained runs are kept as a disclosed,
    gaming-flagged secondary row.
\end{itemize}

\subsection{Selection}
\label{sec:selection}

% src: WT/benchmarks/learned_routing/remote/p2orch.py select (rule string); CR/facts/gate.json
%      full_plan.objective_cells_repeats.selection; CR/CONTRACT.md A17.2; CR/literature/LESSONS.md LR-10,
%      LR-11 action 3; train.py _write_best (best.json keeps best_ever_train and selected_by_val)
Selection happens at two levels, and neither uses the best train sample, whose score is biased
upward by the search itself by an amount comparable to the differences being measured
\citep{cawley2010overfitting}.
\begin{enumerate}
  \item \textbf{Within a policy.} A policy's configuration is the validated candidate with the
    highest validation objective ($\krep = 0,1,2$) over its three restarts: two candidates at each
    of five checkpoints per restart, up to 30 per policy. \code{best.json} keeps the best train
    sample beside the selection, for reference only.
  \item \textbf{Across policies.} The best baseline $\pibase$ is the tuned baseline whose selected
    configuration has the highest mean clipped log-ratio against $\piref$ on \emph{fresh} validation
    replicates $\krep = 3,\dots,10$. The headline learned policy $\pilearn$ is chosen by the same
    rule among the router-observable learned arms (M1, M2, M1-noaff, M1-v2 and any later learned
    arm; Amendment A17.2). No within-policy selection reads replicates 3 to 10, so the cross-policy
    choice is not made on the scores that picked each configuration.
\end{enumerate}
Every learned variant scored on validation is counted, and the number of learned arms compared is
reported with the headline (LR-11).

\subsection{Equal budget}
\label{sec:budget}

% src: CR/facts/gate.json full_plan.budget.equal_footing; WT/benchmarks/learned_routing/remote/p2orch.py
%      budget_report (tier_A_equal_budget: every finished tier-A run has (fevals, val_checkpoints) =
%      (400, 5); the re-evaluation count is reported but excluded); CR/runs/phase2/MANIFEST.json args
Every tuned policy, learned or baseline, gets the same resources: $\budget = 400$ objective
evaluations per restart on all 34 train cells with $\Krep = 2$ replicates, $\popsize = 16$, three
restarts, five validation checkpoints of two candidates on 14 cells with $\Kval = 3$, the same
reference and objective, and the same selection rule. The orchestrator verifies that every finished
run spent exactly 400 evaluations and five validation checkpoints; the re-evaluation diagnostic is
outside this check because a wall-clock chunk can skip it. No run starts from the pilot's tuned
configurations, which would give M0 and M1 the pilot's 480 evaluations on top of $\budget$.

Equal budget fixes what each policy may spend, not how far it can get. The spaces differ in
dimension (2 to 22) and in how strongly decisions respond to the parameters (\cref{tab:spaces}),
and only the learned arms start from points derived from another policy (\cref{sec:restarts}).
Which method wins can depend on the budget \citep{dodge2019show}, so two checks ask whether any
policy was still improving at $\budget$. No tuned baseline of tiers B, C or the \ais{} league was
selected at its last validated generation; the learned arms selected there had gained 0.0001 to
0.0033 over the last validation interval, at most 0.09 of the MDE. And an auditor's validation-oracle
probes, which deliberately favor the baselines, found no headroom: 38 \policy{ramjet} configurations
around all three restarts' bases reached at most 0.1421 against the tuned 0.1420, and widened probes
of M0 and \policy{llm-d-precise-prefix} stayed below their tuned values (0.1299 against 0.1334 and
0.1202 against 0.1414), all far below M1-v2's 0.1767.
% src: CR/facts/tierBC.json decisions.improving_at_cutoff_i9 (0.0001-0.0033, at most 0.086 x MDE);
%      CR/audits/final/refute-headline-2.md ("baselines not under-tuned (val-oracle ceilings: ramjet 38
%      configs incl. all 3 bases max 0.1421 vs tuned 0.1420; M0 20 widened configs max 0.1299 vs 0.1334;
%      llmdpp 12 widened max 0.1202 vs 0.1414; all << M1-v2 0.1767)"), via CR/facts/STATE.md

\subsection{Drift check before M2}
\label{sec:drift}

% src: CR/runs/phase2/MANIFEST.json drift-{l1,l3,n4,n8} (96 evals, --val-every 3, cell/val filters);
%      CR/facts/DEVIATIONS.md 2026-10-03 "LR-15 drift runs use 96 evaluations"; p2orch.py drift_matrix
%      (val k3-10, loss = matrix[b][b] - matrix[a][b], within_mde rule, mde 0.038);
%      train subset sizes from CR/cells/train.jsonl (L1 6, L3 8, N4 18, N8 16 cells);
%      CR/facts/gate.json full_plan.cut_order_if_over[4]; CR/CONTRACT.md A17.3; LESSONS.md LR-15
A context term can help only if the best coefficients change with the regime (LR-15). Before M2,
M1 is tuned at a quarter of the budget on each of four train subsets: L1 cells only (6 cells), L3
only (8), $\Nw = 4$ only (18) and $\Nw = 8$ only (16). The quarter budget is 96 evaluations, six
whole generations, because 100 is not a multiple of 16. Each run validates every 3 generations on
the matching validation cells and selects on $\krep = 0,1,2$. Each selected model is then scored
on every regime's validation cells with fresh replicates $\krep = 3,\dots,10$. The cross-play loss
of the model tuned on regime $a$, evaluated on regime $b$, is the mean clipped log-ratio against
$\piref$ of $b$'s own model on $b$'s validation cells minus that of $a$'s model; the pairs are L1
and L3, and $\Nw = 4$ and $\Nw = 8$. The gate planned to deprioritize M2 if all four losses lay
within the MDE of 0.038. Amendment A17 later made M2 at rank 2 mandatory, so the check no longer
decides whether M2 trains; it is reported as evidence for or against context terms. It came out
inconclusive, with no evidence of drift (\cref{sec:context}).

\subsection{Compute and numerics}
\label{sec:compute}

% src: CR/CONTRACT.md A9, A10; CR/facts/phase2_prep.json targets_at_start (10 remote nodes plus the
%      workstation); CR/facts/STATE.md phase2-prep (early generation walls on EPYC 9654P and EPYC 7702P
%      nodes), tier-B/C relays i6-i9 (about 300 s and about 610-690 s per generation; ingest counts),
%      phase2-mid fix (constrained M1 re-runs 4 h 24 min to 4 h 31 min); CR/audits/phase2-mid/fix.md
Replays run on CPUs: locally through a pool of 20 slots, and on a pool of shared CPU nodes, one
tuning run per node; the first wave started on 10 remote nodes and the workstation. A generation
took 225 to \SI{333}{s} early in phase~2, and about \SI{300}{s} later, on a 96-core EPYC 9654P node,
and 610 to \SI{690}{s} on a 64-core EPYC 7702P node; the three sign-constrained M1 re-runs took
4\,h~24\,min to 4\,h~31\,min each on the latter. Phase 2 ran 66 restarts of 400 evaluations in
total (tiers A and B, the ablation and the \ais{} league), and each run whose ingest the stage log
records stored between \num{30199} and \num{30792} result records. The single test pass added
\num{5220} records on the tuning build and 900 on the \ais{} build, and the robustness pass
\num{23940} (\cref{sec:res-robustness}). Replay-seconds for phase 2 were not totaled.
% src: CR/facts/test_results.json execution.evaluated_once (5,220 + 900); CR/facts/robustness.json
%      execution.records (lag 10,800 + lag-0 900 + timing 4 x 3,060 = 23,940; coverage[*].records sum to the same)

% src: CR/facts/remote.json summary, parity.totals (6,312 records, 48,901,384 rows), caveats[3]
%      (numerics host-dependent; pins equal workstation defaults; 60/60 generations),
%      parity.lr_train.numerics_check (EPYC 9654P default: first diff at generation index 0; pinned,
%      and EPYC 7702P: 60/60 identical; fake objective, budget 720, popsize 12), parity.lr_train.pinned_on_epyc_9654p
%      (best.json, best_policy.yaml, identity.json identical; history identical except gen_wall_s);
%      environment.lr_train_numerics_pin; WT/benchmarks/learned_routing/remote/node_train.sh:73-76
Runs move between machines, so a run must continue bit-exactly on any host. Two things depend on
the host. Replay results depend on the operating system's math library, so every distinct node
image was checked against the workstation before use: across \num{6312} remote records, every field
matched, and \num{48901384} per-request rows matched byte for byte. CMA-ES's own linear algebra
depends on numpy's and OpenBLAS's CPU dispatch. In a 60-generation \code{lr-train} run on a
synthetic objective without replays, an EPYC 9654P node with default settings diverged from the
workstation in the last bits from the first generation, and matched it in all 60 generations once
three environment variables pinned OpenBLAS to a single thread and core type and disabled numpy's
AVX-512 code paths; the pinned values equal the workstation's defaults. With the pins, two short
real \code{lr-train} runs on that node reproduced the workstation's history (apart from wall-clock
fields), \code{best.json} and \code{best\_policy.yaml} exactly.

% src: train.py module docstring ("checkpoint.pkl is written twice per generation ... stores the pycma
%      object and numpy's global RNG state"), Trainer._load_state (identity check); CR/facts/STATE.md
%      audit(build/harness-robustness-splits r0) ("lr-train kill -9 x3 resume bit-identical (history,
%      best, CMA state, continuation)")
\code{lr-train} writes its checkpoint twice per generation, after sampling and after the update,
including the CMA-ES object and numpy's random state, so a resumed run continues exactly as an
uninterrupted one would; an audit killed a run three times and found its history, selection,
CMA-ES state and continuation identical. A run refuses to resume if its space, cells, reference,
objective or seeds have changed.

% src: CR/facts/gate.json full_plan.cut_order_if_over[6]; CR/facts/pilot.json
%      measurements.subset_vs_full_objective (spearman 0.947, 16 candidates)
If compute ran short, the last resort in the gate's cut order is to score each generation on a
rotating 12-cell subset of the train cells, for every policy of a wave alike. On the 16 candidates of
one pilot generation, subset and full-train objectives had a Spearman correlation of 0.947
(\emph{preliminary (pilot, train subset)}).
